\section{Conclusion}

This thesis has presented a comprehensive four-objective optimization framework
for the strategic deployment of Electric Vehicle Charging Stations across the
85~administrative wards of Indore, Madhya Pradesh. The framework unifies
logistic EV adoption forecasting, 100-scenario Monte~Carlo stochastic demand
simulation, M/M/$c$/K finite-capacity queueing analysis with Time-of-Use
demand segmentation, and two state-of-the-art multi-objective evolutionary
algorithms --- NSGA-II and MOPSO --- into a single, deployment-ready
computational pipeline.

The four-objective formulation --- minimizing total capital expenditure and
mean queue waiting time while maximizing spatial demand coverage and net
return on investment --- successfully captures the conflicting trade-offs
inherent in real-world EVCS planning. Replacing raw annual profit with net
ROI (annual profit / CAPEX) as the financial objective prevented Pareto
front degeneracy and produced genuine solution diversity across the
cost--coverage--profitability--service space. Embedding the M/M/$c$/K
blocking probability constraint ($P_{\text{block}} \leq 10\%$) directly
inside the evolutionary evaluation loop, rather than applying it
post-hoc, proved critical: it forced both algorithms to provision
adequate charger capacity at demand bottlenecks, yielding
operationally stable deployment plans under peak load conditions.

Formal performance evaluation using Hypervolume, Generational
Distance, and Spread indicators showed that NSGA-II and MOPSO are
\emph{complementary} rather than dominant on this binary
combinatorial problem. MOPSO achieves a marginally higher final
normalised hypervolume of $1.0434$ versus $0.9967$ for NSGA-II,
and a substantially lower generational distance to the union
reference front. NSGA-II, in turn, produces a more uniformly
spread front (Spread $0.750$ vs.\ $0.792$) and a slightly lower
median queue waiting time ($1.45$~min vs.\ $1.79$~min). Across the
union of the two archives, capital expenditures range from
\rupee39\,L to \rupee411\,L, coverage from 46.4\,\% to 100\,\%,
and Net~ROI from 1.85 to 4.33. All Pareto-optimal layouts satisfy
the $P_{\text{block}} \leq 10\,\%$ blocking constraint and the
$u < 1$ stability bound across the three Time-of-Use periods,
confirming that embedding the M/M/$c$/K constraint inside the
evaluation loop successfully steers the search away from
queue-infeasible deployments. The annual-profit distribution
across the union archive spans \rupee1.2\,Cr to \rupee11.6\,Cr per
year, confirming the commercial viability of public-private
partnership EVCS deployment in Tier-2 Indian cities.

Sensitivity analysis confirmed robustness to EV growth rate, electricity
tariff, and charger technology variations, while scalability benchmarks
demonstrated polynomial runtime scaling for both algorithms, validating
the framework's applicability to larger metropolitan deployments.
Overall, this study demonstrates that rigorous multi-objective evolutionary
optimization, combined with realistic stochastic demand modelling and
analytical queueing constraints, can produce actionable, evidence-based
EVCS deployment strategies that simultaneously satisfy spatial, operational,
and financial requirements.

%---------------------------------------------------------
\section{Future Work}

Several extensions can further enhance the scope, accuracy, and practical
utility of the proposed framework:

\begin{itemize}
    \item \textbf{Road-network routing:} The current framework computes
          spatial coverage using Euclidean distances. Integrating
          OpenStreetMap-based road-network routing via graph-search
          APIs (e.g., OSRM) would produce topologically accurate
          service radii, particularly in areas with low road density
          or irregular urban form.

    \item \textbf{Battery swapping station integration:} For the Indian
          two-wheeler and three-wheeler market --- the dominant EV
          segments --- expanding the optimization scope to include
          Battery Swapping Stations alongside conventional plug-in
          chargers would significantly improve the framework's
          coverage of actual urban charging demand.
    \item \textbf{Dynamic and real-time pricing:} The current ToU pricing
          structure is applied exogenously based on time-of-day patterns.
          Future work could incorporate real-time dynamic pricing
          responsive to localized substation load levels and live
          demand, potentially modelled through a Stackelberg game
          between the network operator and EV users.

    \item \textbf{Research publication:} Based on the results presented
          in this thesis, a full research paper will be prepared and
          submitted to a peer-reviewed journal or international
          conference in smart cities, transportation electrification,
          or evolutionary computation, documenting the four-objective
          formulation, NSGA-II vs.\ MOPSO comparison, and Indore
          case study findings.
\end{itemize}

\noindent
These extensions can transform the current framework into a full-scale
EV infrastructure intelligence platform, supporting sustainable,
data-driven mobility planning for the rapidly electrifying urban
ecosystems of Indian and other emerging-market cities.